\documentclass[11pt]{article}
\input{style-pc}

\begin{document}

\entetesujet{Sujet bac 2018 -- Physique-Chimie -- Série C}

% ============================ CHIMIE ============================
\matiere{CHIMIE}{8 points}

\partie{A : vérification des connaissances}

\noindent\textbf{1. Questions à alternative vrai ou faux.} Réponds par vrai ou faux aux affirmations suivantes. Exemple : 1.f = Vrai.
\begin{enumerate}[label=1.\alph*)]
  \item Une réduction est une réaction au cours de laquelle il y a un gain d'électrons.
  \item Les lois de Raoult ne s'appliquent qu'aux solutions diluées non électrolysables.
  \item Pour un mélange équimolaire, le rendement d'une réaction d'hydrolyse vaut 33 \% lorsqu'il s'agit d'un alcool secondaire.
  \item Pour une réaction d'ordre un, le temps de demi-réaction est $\dfrac{1}{kC_0}$.
\end{enumerate}

\medskip
\noindent\textbf{2. Texte à trous.} Recopie puis complète le texte ci-après par quatre des mots manquants suivants : \textit{ionisation ; énergie ; transitions ; raies ; excité ; radiations}.

\textit{L'ensemble des $\cdots\cdots$ émises lors des $\cdots\cdots$ aboutissant au même niveau d'$\cdots\cdots$ constitue une série de $\cdots\cdots$.}

\medskip
\noindent\textbf{3. Question à réponse courte.} Donne la relation qui lie le pH au pKa d'un couple acide/base.

\partie{B : application des connaissances}
\noindent Données : masses molaires atomiques (en g$\cdot$mol$^{-1}$). Mn : 55 ; K : 39 ; O : 16.

On prépare une solution de permanganate de potassium $(\ch{K^+} + \ch{MnO_4^-})$ en dissolvant une masse $m = 19{,}75$ g de cristaux anhydres de permanganate de potassium ($\ch{KMnO_4}$) dans 250 mL d'eau distillée.

\begin{enumerate}
  \item Calcule la concentration $C_1$ de la solution ainsi obtenue.
  \item On veut déterminer la concentration d'une solution d'eau oxygénée ($\ch{H_2O_2}$). Pour ce faire, on prélève un volume $V_2 = 10{,}0$ mL de cette solution que l'on fait réagir avec la solution de permanganate de potassium. Le volume total de solution de permanganate de potassium versé à l'équivalence est $V_1 = 8{,}0$ mL. Les couples redox mis en jeu sont : $\ch{MnO_4^-/Mn^{2+}}$ et $\ch{O_2/H_2O_2}$.
  \begin{enumerate}[label=\alph*.]
    \item Écris les demi-équations électroniques intervenant.
    \item Déduis l'équation-bilan de la réaction d'oxydoréduction qui a lieu.
    \item Détermine la concentration $C_2$ de la solution d'eau oxygénée.
  \end{enumerate}
\end{enumerate}

% ============================ PHYSIQUE ============================
\matiere{PHYSIQUE}{12 points}

\partie{A : vérification des connaissances}

\noindent\textbf{1. Questions à choix multiples.} Choisis la bonne réponse parmi les propositions suivantes. Exemple : 1.5 = a.
\begin{enumerate}[label=1.\arabic*)]
  \item La pulsation propre d'un pendule de torsion est :
  \begin{enumerate}[label=\alph*)]
    \item $\omega_0 = \sqrt{\dfrac{g}{l}}$ ;
    \item $\omega_0 = \sqrt{\dfrac{k}{m}}$ ;
    \item $\omega_0 = \sqrt{\dfrac{c}{J_\Delta}}$.
  \end{enumerate}
  \item Un circuit RLC est dit résonnant si :
  \begin{enumerate}[label=\alph*)]
    \item $L\omega - \dfrac{1}{c\omega} > 0$ ;
    \item $L\omega - \dfrac{1}{c\omega} < 0$ ;
    \item $L\omega = \dfrac{1}{c\omega}$.
  \end{enumerate}
  \item Dans une région où règne un champ électrique uniforme $\vec{E}$, à l'extérieur des plaques le vecteur champ électrique $\vec{E}$ est :
  \begin{enumerate}[label=\alph*)]
    \item inférieur à zéro ;
    \item supérieur à zéro ;
    \item nul.
  \end{enumerate}
  \item Dans un mouvement rectiligne uniformément accéléré, le produit scalaire $\vec{v}\cdot\vec{a}$ est :
  \begin{enumerate}[label=\alph*)]
    \item nul ;
    \item inférieur à zéro ;
    \item supérieur à zéro.
  \end{enumerate}
\end{enumerate}

\medskip
\noindent\textbf{2. Texte à trous.} Complète la phrase suivante par quatre des cinq mots suivants : \textit{mécanique ; cinétique ; conservatif ; constante ; temps}.

\textit{Un système est dit $\cdots\cdots$ lorsque son énergie $\cdots\cdots$ reste $\cdots\cdots$ au cours du $\cdots\cdots$}

\partie{B : application des connaissances}
Une lame vibrante est animée d'un mouvement sinusoïdal de fréquence $f$, de période $T$ et d'amplitude $a$. Elle est munie d'une pointe qui frappe verticalement la surface d'une nappe d'eau en un point $S$. On suppose qu'il n'y a ni amortissement, ni réflexion des ondes. À la date $t = 0$, le point $S$ commence à vibrer. Ces oscillations se propagent à la surface de l'eau avec la célérité $c = 20$ m/s. Les graphes ci-dessous représentent l'état vibratoire des points $S$ et $M$. Le point $M$ est situé à distance $\overline{SM} = x$ de la source $S$.

\begin{center}
\begin{tikzpicture}[>=Stealth,scale=1]
  % axes
  \draw[->] (-0.3,0)--(9.4,0) node[right]{$t$};
  \draw[->] (0,-1.9)--(0,1.9) node[above]{$y$(mm)};
  \draw[dashed] (0,1.3)--(9.2,1.3);
  \draw[dashed] (0,-1.3)--(9.2,-1.3);
  \node[left] at (0,1.3){$2$};
  \node[left] at (0,0){$0$};
  \node[left] at (0,-1.3){$-2$};
  % yS solide : 2 sin(pi x), periode 2
  \draw[thick] plot[domain=0:9,samples=160] (\x,{1.3*sin(deg(pi*\x))});
  % yM dashed : retardé de 0.8
  \draw[thick,dashed] plot[domain=0.8:9,samples=150] (\x,{1.3*sin(deg(pi*(\x-0.8)))});
  % labels
  \node at (1.4,1.65){$y_S(t)$};
  \node at (3.4,1.65){$y_M(t)$};
  % periode 0.015 s
  \draw[<->] (0,-1.75)--(2,-1.75) node[midway,below]{$0{,}015$ s};
\end{tikzpicture}
\end{center}

\begin{enumerate}
  \item À partir de la courbe $y_S(t)$ :
  \begin{enumerate}[label=\alph*.]
    \item Détermine les valeurs de la période $T$ et de l'amplitude $a$.
    \item Avec quel retard $\theta$ par rapport à $S$, le point $M$ commence-t-il à vibrer ?
    \item D'après les courbes $y_S(t)$ et $y_M(t)$, quel déphasage existe-t-il entre les mouvements de $S$ et de $M$ ?
    \item Établis les équations horaires $y_S(t)$ et $y_M(t)$.
  \end{enumerate}
\end{enumerate}

\partie{C : résolution du problème}
On se propose de déterminer la longueur d'onde $\lambda_2$ inconnue d'une radiation monochromatique. Pour cela, on réalise l'expérience d'interférences lumineuses au moyen du dispositif de fentes d'Young.

\begin{center}
\begin{tikzpicture}[>=Stealth,scale=1]
  \draw[dashed] (-0.3,0)--(7.6,0);
  \fill (0,0) circle (1.5pt) node[left]{$S$};
  \draw[thick] (3,-1.6)--(3,-0.5);
  \draw[thick] (3,0.5)--(3,1.6);
  \node at (2.7,0.9){$S_1$};
  \node at (2.7,-0.9){$S_2$};
  \draw[<->] (3.25,-0.45)--(3.25,0.45) node[midway,right]{$a$};
  \draw[thick] (7.4,-1.8)--(7.4,1.8) node[above right]{$(E)$};
  \draw[<->] (0,-1.4)--(3,-1.4) node[midway,below]{$d$};
  \draw[<->] (3,-1.4)--(7.4,-1.4) node[midway,below]{$D$};
\end{tikzpicture}
\end{center}

\noindent La distance séparant les fentes $S_1$ et $S_2$ est $a = 1$ mm. L'écran $(E)$ d'observation est placé à la distance $D = 2$ m du plan des fentes.

\begin{enumerate}
  \item La source $S$ émet d'abord une radiation lumineuse monochromatique de longueur d'onde $\lambda_1$. On mesure la distance séparant 11 franges brillantes consécutives, on trouve $l = 9{,}6\cdot10^{-3}$ m.
  \begin{enumerate}[label=1.1.]
    \item \begin{enumerate}[label=\alph*.]
      \item Écris la formule donnant la position des franges brillantes sur l'écran en fonction de $\lambda_1$, $D$ et $a$.
      \item Établis l'expression de l'interfrange $i$ en fonction de $\lambda_1$, $D$ et $a$.
    \end{enumerate}
  \end{enumerate}
  \begin{enumerate}[label=1.2.]
    \item \begin{enumerate}[label=\alph*.]
      \item Calcule la valeur de l'interfrange.
      \item Déduis la valeur de la longueur d'onde $\lambda_1$ de la radiation émise.
    \end{enumerate}
  \end{enumerate}
  \item La source $S$ émet simultanément les deux radiations lumineuses de longueurs d'onde $\lambda_1$ et $\lambda_2$. Les deux systèmes de franges se superposent sur l'écran. On constate que la cinquième (5\textsuperscript{e}) frange brillante de la radiation de longueur d'onde $\lambda_1$ et la quatrième (4\textsuperscript{e}) frange brillante de la radiation de longueur d'onde $\lambda_2$ coïncident. Détermine la longueur d'onde $\lambda_2$.
\end{enumerate}
\noindent Donnée : 1 mm $= 10^{-3}$ m.

\end{document}
